% TOP_PROBLEM: {"id": 4700022, "title": "Triangular billiards", "category_id": 11, "category": {"id": 11, "name": "dynamical_systems", "display_name": "Dynamical Systems", "description": "Problems about long-term behavior of deterministic systems, Hamiltonian dynamics, and periodic orbits.", "slug": "dynamical-systems", "order_index": 11, "created_at": "2026-07-31T15:26:25.671Z"}, "status": "partially_solved", "problem_number": "AMR-046-0022", "set_id": 13}
% FIRST_POSTED: 2026-10-02
% ENTRY_KIND: statement_only
% UnsolvedMath ID 4700022; source catalogue code AMR-046-0022
% !TeX program = lualatex
% Definition and sources only; this file is not an LLM solution attempt.
\documentclass[11pt,a4paper]{article}
\usepackage[margin=25mm]{geometry}
\usepackage{fontspec}
\setmainfont{lmroman10-regular.otf}[BoldFont=lmroman10-bold.otf,
 ItalicFont=lmroman10-italic.otf,BoldItalicFont=lmroman10-bolditalic.otf]
\usepackage{amsmath,amssymb,mathtools}
\usepackage{unicode-math}
\setmathfont{latinmodern-math.otf}
\AtBeginDocument{\let\setminus\smallsetminus}
\usepackage{xurl,hyperref}
\hypersetup{colorlinks=true,urlcolor=blue}
\setcounter{secnumdepth}{0}
\setlength{\parindent}{0pt}
\setlength{\parskip}{5pt}
\setlength{\emergencystretch}{3em}
\begin{document}
\section*{Triangular billiards}\label{4700022}
{\small UnsolvedMath ID 4700022 · AMR-046-0022 · Dynamical Systems · AMR Open Problem Lists}
\subsection{Definitions and mathematical statement}
Let $\Delta\subset\mathbb R^2$ be a nondegenerate closed Euclidean triangle.
A billiard trajectory moves at constant positive speed on straight
segments in its interior. At a collision with the relative interior of a
side it reflects specularly: the tangential component of velocity is
preserved and the normal component changes sign. Trajectories that meet
a vertex are excluded, because the reflection rule is then undefined.

A periodic trajectory is a finite sequence of nonzero segments that
returns to its starting position with its starting velocity and thereafter
repeats. Equivalently, there are a starting state $(x,v)$ and $T>0$ for
which the nonsingular billiard flow satisfies
\[
 \Phi_T(x,v)=(x,v).
\]
The question is whether every nondegenerate triangle admits at least one
such periodic trajectory. The initial position and direction may depend
on the triangle; the assertion is not that every trajectory is periodic.

No rationality condition is imposed on the angles, and the number of
reflections is allowed to be arbitrarily large. Self-intersections of the
geometric path are allowed, while corner collisions and zero-length
segments are not. The classical orthic construction handles acute
triangles, and other classes have known positive results. The target
includes arbitrary obtuse triangles. A near-return found numerically is
not enough: one must verify an exact closed orbit that remains away from
all three vertices throughout its reflections.
\subsection{Short English statement}
Does every triangle admit a closed billiard path that avoids its vertices and returns with the same direction?
\subsection{Sources}
\begin{itemize}
\item[S1] ULAM.AI and the UnsolvedMath Contributors, supplied problems.json, record 4700022 (AMR-046-0022), AMR Open Problem Lists. Catalogue identity and source transcription; CC BY 4.0.
\url{https://huggingface.co/datasets/ulamai/UnsolvedMath}.
\item[S2] Gasull - Some open problems in low dimensional dynamical systems (2020), Problem environment 22: Triangular billiards. Original source indexed by AMR-046-0022.
\url{https://arxiv.org/abs/2012.02524}.
\end{itemize}
\end{document}
